\section{Lean formalization architecture}\label{sec:lean}

\subsection{Pinned upstream API}

The outline targets the following namespaces and theorem interfaces in the audited snapshot:
\begin{itemize}
  \item \texttt{LeanProofs.ConciseIPRA.PR}, \texttt{PTm}, \texttt{PRAAx}, and \texttt{PRA.qfInd};
  \item \texttt{LeanProofs.ConciseIPRA.Prf.subst}, \texttt{Prf.trans}, equality and logical constructors;
  \item \texttt{LeanProofs.IPRA.inst} and \texttt{LeanProofs.IPRA.ind};
  \item \texttt{LeanProofs.IPRA.Search.mu}, \texttt{search\_spec}, \texttt{search\_found}, and \texttt{search\_minimal};
  \item the pairing, bounded-sum, order, quotient, integer, and descent libraries;
  \item \texttt{LeanProofs.IPRA.closeAll}, \texttt{deriveNat}, and \texttt{ipra\_ring};
  \item \texttt{pra\_iff\_native} for the final transfer.
\end{itemize}
The formalization should live beside \RepoPath{Derivations/Lean/IPRADerivations}, not copy the syntax or deduction kernel.

\subsection{Module dependency graph}

A concrete implementation order is:
\begin{enumerate}[label=\textbf{M\arabic*.}]
  \item \texttt{KIFC/KernelCore.lean}: local rules, proof schedules, implication contracts, loop certificates.
  \item \texttt{KIFC/QFBoolean.lean}: Boolean symbols, range/reflection, and syntax-directed compilation of fixed formulas.
  \item \texttt{KIFC/BoundedFold.lean}: \texttt{allB}, \texttt{anyB}, table folds, and reflection.
  \item \path{KIFC/BoundedExists.lean}: truth-to-zero adapters and wrappers for the checked search theorems.
  \item \texttt{KIFC/FiniteTable.lean}: fixed tuple/table coding and projection/write equations.
  \item \texttt{KIFC/Trace.lean}: generic packed-state trace theorem using \texttt{IPRA.ind}.
  \item \texttt{KIFC/Demand.lean}: reverse demand on a fixed generated DAG.
  \item \texttt{KIFC/FiniteRelation.lean}: viability rows, path checker, and least viable path.
  \item \texttt{KIFC/Observer.lean}: exact stability, diameter stability, guarded labels, and finite gluing.
  \item \texttt{KIFC/RationalMetric.lean}: signed rational codes, triangle inequalities, finite sums, tail budgets.
  \item \texttt{KIFC/RetainedTransfer.lean}: negative stability, proof translation, recovery, and retained transfer.
  \item \texttt{KIFC/Compiler.lean}: typed metalevel graph and specialized proof-producing elaborator.
  \item theorem-specific generated modules and \texttt{Checks/Axioms.lean}.
\end{enumerate}
The packet includes source-oriented artifacts for the kernel, Boolean/search adapter, bounded-existential, relational, graph, retained-transfer, and example layers.  \Cref{sec:repo-integration} maps these modules to the exact current repository services and gives exit criteria.

\subsection{Kernel rules and schedules}

The smallest useful object is a local rule.

\begin{lstlisting}
structure KernelRule (n : Nat) where
  premises   : List (Fm PTm n)
  conclusion : Fm PTm n
  sound      : premises |- [PRAAx] conclusion

def KernelRule.use (r : KernelRule n)
    (supply : forall psi, psi in r.premises -> Gamma |- [PRAAx] psi) :
    Gamma |- [PRAAx] r.conclusion :=
  r.sound.trans supply
\end{lstlisting}
(The ASCII turnstile above abbreviates the repository notation.)  This wrapper makes the proof boundary visible in types.

A source-complete schedule can be represented by an inductive family:

\begin{lstlisting}
inductive Schedule (n : Nat) (Gamma : List (Fm PTm n)) :
    Fm PTm n -> Type
  | done  (h : Gamma |- [PRAAx] phi) : Schedule n Gamma phi
  | have_ (hpsi : Gamma |- [PRAAx] psi)
      (rest : Schedule n (psi :: Gamma) phi) : Schedule n Gamma phi

noncomputable def Schedule.compile : Schedule n Gamma phi ->
    Gamma |- [PRAAx] phi
  | .done h => h
  | .have_ hpsi rest =>
      (Schedule.compile rest).trans (fun chi hchi => by
        cases hchi with
        | head => exact hpsi
        | tail h => exact Prf.hyp h)
\end{lstlisting}
The actual file uses the exact list-membership constructors available in the pinned Lean version.  This theorem should be compiled before any graph AST is attempted.

\subsection{Loop certificates}

The current derivation library already provides exactly the required induction wrapper.  A reusable structure is:

\begin{lstlisting}
structure LoopCertificate (k : Nat) where
  inv  : Fm PTm (k + 1)
  qf   : inv.QF
  zero : ([] : List (Fm PTm k)) |- [PRAAx]
           inv.subst (Syn.inst PTm.zero)
  succ : [inv] |- [PRAAx]
           inv.subst (Syn.repl (PTm.succ (Syn.var 0)))

def LoopCertificate.compile (c : LoopCertificate k) :
    ([] : List (Fm PTm (k + 1))) |- [PRAAx] c.inv :=
  IPRA.ind c.inv c.qf c.zero c.succ
\end{lstlisting}
Instantiation at an arbitrary time and parameters is then \texttt{IPRA.inst c.compile sigma}.  The generated trace theorem should reduce to constructing this structure.

\subsection{Boolean reflection API}

Avoid a first attempt at a fully internal formula-code evaluator.  Start with fixed combinators:

\begin{lstlisting}
structure Reflected (n : Nat) where
  formula : Fm PTm n
  code    : PTm n
  range   : ([] : List (Fm PTm n)) |- [PRAAx]
              (code = 0) or (code = 1)
  iff_one : ([] : List (Fm PTm n)) |- [PRAAx]
              Fm.iff formula (code = 1)
\end{lstlisting}
Implement constructors \texttt{atomEq}, \texttt{not}, \texttt{and}, \texttt{or}, and later bounded folds.  The syntax-directed compiler can remain a metalevel function returning \texttt{Reflected n}; no object formula codes are needed.

For bounded folds, define symbols using \texttt{PR.prec}.  The body predicate is a fixed \texttt{PR} symbol or fixed term template, so the fold symbol is generated per predicate.  Prove its zero/successor equations with \texttt{PRA.precZero} and \texttt{PRA.precSucc}.  Prove Boolean range and reflection once for the generated fold schema.

\subsection{Checked-search adapter and bounded witness certificates}

The audited repository already has the numerical least-zero function \texttt{Search.mu} and object theorems \texttt{search\_spec}, \texttt{search\_found}, and \texttt{search\_minimal}.  KIFC should not maintain a second generic search proof.  Given a reflected Boolean code \(\chi_P\), define the zero-on-truth body
\[
  z_P\defeq 1\mathbin{\dot-}\chi_P.
\]
The Boolean range theorem and monus equations prove \(z_P=0\leftrightarrow\chi_P=1\leftrightarrow P\).  Instantiate the checked search theorems with \(z_P\), then package the transported conclusions:

\begin{lstlisting}
structure SearchCertificate (n : Nat) where
  pre : Fm PTm n
  found witness : PTm n
  hasWitness witnessValid witnessMinimal : Fm PTm n
  foundIff : [pre] |- [PRAAx]
    Fm.iff (found = 1) hasWitness
  valid : [pre, found = 1] |- [PRAAx] witnessValid
  minimal : [pre, found = 1] |- [PRAAx] witnessMinimal
\end{lstlisting}

The bounded-existential wrapper adds a quantifier-free certificate formula recording foundness, the witness bound, reflected validity, and the all-fold asserting no earlier witness.  Its two public equivalences are exactly \cref{thm:bounded-ex-normal}.  Multi-parameter bodies may be generated directly at the needed arity or use the checked pairing projections; no object-level syntax evaluator is involved.

\subsection{Relational viability}

Use a rectangular table code.  The fixed generated parameters are:
\begin{itemize}
  \item horizon term \(H\), maximum row width \(B\);
  \item Boolean validity \(V(i,a)\), edge \(E(i,a,b)\), and terminal \(T(a)\);
  \item table update writing row \(H-j\) at reverse time \(j\).
\end{itemize}
The core QF invariant states that all cells already written equal the corresponding bounded-\texttt{anyB} recurrence and all unwritten cells retain a default.  A semantic list theorem may be proved first, but the public object theorem must separately reflect the table equations.

The target signatures are:

\begin{lstlisting}
theorem viable_iff_suffix :
  [] |- [PRAAx] Fm.iff (viable H i a = 1)
    (boundedPathExists H i a)

theorem least_path_valid :
  [anyStartViable H = 1] |- [PRAAx] pathChecker H (leastPath H) = 1
\end{lstlisting}
The path-existence formula is bounded and can be represented by a coded tuple plus a Boolean checker.  The theorem need not quantify over Lean lists.

\subsection{Observers and metric budgets}

An exact observer record can store the object theorem directly:

\begin{lstlisting}
structure ExactObserver (n : Nat) where
  bound : PTm n
  valid : Fm PTm (n + 1)
  obs   : PTm (n + 1)
  total : [pre] |- [PRAAx] boundedEx bound valid
  stable : validA :: validB :: pre.shift.shift :: [] |- [PRAAx]
             obsA = obsB
\end{lstlisting}
The exact binder/substitution layout should be hidden behind named helper substitutions; do not construct \texttt{Fin.cases} expressions in every application file.  First prove a version where candidate terms \(a,b\) are explicit free variables, then universally close it.

For diameter observers, define signed rational codes and an integer inequality formula.  Reuse \texttt{IPRADerivations.Core.Integers}; prove triangle inequalities and finite-sum monotonicity as object derivations.  The shrinking theorem is a loop over a finite rational sum plus a supplied tail certificate.

\subsection{Graph elaborator}

Do not begin with a dependently typed universal AST.  First generate two or three concrete graph modules.  The metalevel generator should emit:
\begin{enumerate}[label=(\arabic*)]
  \item definitions of specialized \texttt{PR} symbols and \texttt{PTm} projections;
  \item local node formulas;
  \item local proof constructors invoking shared theorems;
  \item a \texttt{LoopCertificate};
  \item a \texttt{Schedule};
  \item the terminal theorem and axiom-audit commands.
\end{enumerate}
After these compile, abstract the repeated code into:

\begin{lstlisting}
inductive NodeKind
  | algebra | boolean | search | loop | relation | observer

structure GraphNode where
  kind : NodeKind
  deps : Array Nat
  -- metalevel descriptions of generated terms and theorem constructors

structure CheckedGraph where
  nodes : Array GraphNode
  topological : Topological nodes
  aritySafe : AritySafe nodes
\end{lstlisting}
The elaborator returns Lean syntax for a specialized module.  Its correctness is checked indirectly because all generated proof terms are elaborated and kernel checked.  A later theorem may verify an internal graph AST, but it is not a prerequisite for useful checked examples.

\subsection{Algebra node adapter}

The current example pattern is:

\begin{lstlisting}
import IPRADerivations.Geometry.TriangleFormulas

open LeanProofs.ConciseIPRA LeanProofs.IPRA

def twoSquaresRule : KernelRule 8 where
  premises := []
  conclusion := Geometry.TriangleFormulas.twoSquaresFormula
  sound := Geometry.TriangleFormulas.twoSquares_open
\end{lstlisting}
New polynomial identities can define their formula and use \texttt{ipra\_ring} in the \texttt{sound} field.  This is the first node class that can immediately inherit a checked implementation.

\subsection{Final native transfer}

After constructing
\[
 h:\Gamma\vdash[\PRAAx]\varphi,
\]
the final native theorem is simply the forward direction of \texttt{pra\_iff\_native} applied to \(h\).  Do not perform the translation earlier: ordinary \iPRA\ has the most mature arithmetic and natural-deduction API, while the native core is best treated as the certified endpoint.

\subsection{Formalization completion criteria}

A module is complete only when:
\begin{enumerate}[label=(\arabic*)]
  \item every public mathematical formula is explicit;
  \item every public theorem concludes \(\Prf\,\PRAAx\) or is clearly marked as a metatheorem about proof construction;
  \item there is no \texttt{sorry}, theorem-specific \texttt{axiom}, or opaque semantic replacement for an object proof;
  \item \texttt{\#print axioms} is recorded;
  \item the module builds from the pinned repository commit;
  \item native transfer occurs after the ordinary derivation;
  \item reference tests are not cited as proof of object soundness.
\end{enumerate}
